
\needspace{4\baselineskip}
\color{uniblue}{\subsection{Контроль наличия напряжения на 1(2) секции шин (КНН 1сш, КНН 2сш)}}
\color{black}

\begin{enumerate}[label=\arabic{section}.\arabic{subsection}.\arabic*, labelsep=4pt, leftmargin=0pt, itemindent=57pt]

\item
КНН 1(2)сш контролирует наличие напряжения на 1(2) секции шин с дополнительным контролем симметрии по $U_2$. В устройстве контроль наличия напряжения на 1(2) секции шин представлен функциями «КНН 1сш» и «КНН 2сш» соответственно. Функции используются для АВР, а также для оперативного включения.

\item
Пусковые условия функций формируются без выдержки времени при превышении всех линейных напряжений на секции шин выше уставки и одновременном выполнении следующих условий, если уставка SGF2 выставлена в положение <<Измеренный>>:
\begin{itemize}
    \item функция введена в работу;
    \item минимальное из подведенных линейных напряжений, измеренных с ТН секции шин, выше уставки;
    \item отсутствует срабатывание ИО напряжения обратной последовательности.
\end{itemize}
\item
Предусмотрена блокировка функций при срабатывании ИО обратной последовательности напряжения выше уставки.

\item
Если уставка SGF2 выставлена в положение <<Внешний>>, работа функции осуществляется при наличии активного внешнего дискретного сигнала <<Наличие U 1(2)сш>> без учета срабатывания измерительных органов.

\item
Ввод функций в работу на этапе параметрирования устройства осуществляется своим программным переключателем SGF1 <<Ввод\_функции>>. Информация о введенной функции отображается выходным сигналом <<КНН 1(2)сш: Ввод>>.

\item
Функции могут быть выведены из работы оперативно путем активации сигнала <<ОВ КНН 1(2)сш>>, что характеризуется активным состоянием сигнала <<КНН 1(2)сш: Оперативный вывод>> на выходе функции.

\item
Алгоритм работы функций выполнен в соответствии с рисунком \ref{knnsh:img1}. В таблице \ref{knnsh:tbl1} приведены параметры, необходимые для настройки функций.

\vspace{3mm}
\begin{figure}[H]
\centering
\includegraphics[width=0.9\textwidth,height=0.9\textheight,keepaspectratio]{\fbpath/407. КНН сш/_latex/img/knnsh.pdf}
\captionof{figure}{Функциональная схема логики функции <<КНН 1(2)сш>>}\label{knnsh:img1}
\end{figure}

\small
\begin{longtable}{|>{\centering\arraybackslash}m{5.3cm}|>{\centering\arraybackslash}m{3.3cm}|>{\centering\arraybackslash}m{4.2cm}|>{\centering\arraybackslash}m{1.8cm}|>{\centering\arraybackslash}m{1cm}|}
\caption{Параметры для настройки функции <<КНН 1(2)сш>>\hfill\vspace{-0.5\baselineskip}}\label{knnsh:tbl1}\\ 
\hline
\rowcolor{gray!30}
Параметр (Параметр на ИЧМ) & Условное обозначение на схеме & Значение/ Диапазон & Единица измерения & Шаг \\ 
\hline
\endfirsthead
\caption*{\hspace{3pt}\emph{Продолжение таблицы \ref{knnsh:tbl1}\hfill\vspace{-0.5\baselineskip}}} \\ % сделано по ГОСТ 2.105 п.6.8.7
\hline
\rowcolor{gray!30}
Параметр (Параметр на ИЧМ) & Условное обозначение на схеме & Значение/ Диапазон & Единица измерения & Шаг \\ 
%\hline
\endhead
%\hline
\endfoot
\hline
\endlastfoot
%==+t1*PTOV1|SV_VOLPRES1
\centering Ввод функции в работу (Ввод\_функции)  & \centering SGF1 & \centering Не предусмотрено \\ Предусмотрено & \centering -- & \centering \arraybackslash -- \\
\hline
\centering Напряжение срабатывания (Uср)  & \centering U> & \centering 50,0...150,0 & \centering В & \centering \arraybackslash 0,1 \\
\hline
\centering Напряжение срабатывания ОП (U2ср)  & \centering U2> & \centering 2,0...50,0 & \centering В & \centering \arraybackslash 0,1 \\
\hline
\centering Режим пуска (Реж\_пуска)  & \centering SGF2 & \centering Измеренный \\ Внешний & \centering -- & \centering \arraybackslash -- \\
%===t1
\end{longtable}
\normalsize

\end{enumerate}
